
\renewcommand{\Folder}{2025_RubikCode_Grafo}
\renewcommand{\EntradaBibtex}{CuboRubik_Grafo}

\begin{frame}{\citetitle{\EntradaBibtex} \footnotemark[1] (1)}
%\begin{block}{Brazo Robótico \footnotemark} 
\begin{columns}
	\begin{column}{0.55\textwidth}
		\begin{itemize}
			\item Se modeló un cubo de Rubik cuyo estado puede ser cambiado por el usuario
			\item Se generó un grafo que representa fielmente el estado del cubo de las caras del cubo de rubik
		\end{itemize}
	\end{column}
	\begin{column}{0.44\textwidth}
		 \begin{center}
		     \includegraphics[width=0.98\textwidth]{\Folder/figs/cuadro}
		 \end{center}
	%\end{column}
	%\begin{column}{0.22\textwidth}  
		%\begin{center}
		 %\includegraphics[width=0.98\textwidth]{\Folder/figs/cuadroG}
		 %\end{center}
	\end{column}
\end{columns}
%\end{block} 
\footnotetext[1]{\fullcite{\EntradaBibtex}}
\end{frame}

\begin{frame}{\citetitle{\EntradaBibtex} \footnotemark[1] (2)}
%\begin{block}{Brazo Robótico \footnotemark} 
\begin{columns}
	\begin{column}{0.55\textwidth}
		\begin{itemize}
			\item Los movimentos de las caras del cubo de Rubik son animados para mejor entendimiento del usuairo
			\item Conforme el usuario manipula las caras del cubo de Rubik, el grafo es actualizado
		\end{itemize}
	\end{column}
	\begin{column}{0.44\textwidth}
		 %\begin{center}
		 %    \includegraphics[width=0.98\textwidth]{\Folder/figs/cuadro}
		 %\end{center}
	%\end{column}
	%\begin{column}{0.22\textwidth}  
		\begin{center}
		 \includegraphics[width=0.98\textwidth]{\Folder/figs/cuadroG}
		 \end{center}
	\end{column}
\end{columns}
%\end{block} 

\end{frame}

